\chapter{Referencial Teórico}
\label{cap:referencial}

Este capítulo apresenta os dois algoritmos de ordenação objeto deste trabalho ---
o Insertion Sort e o Quicksort --- e a análise de complexidade de ambos, com ênfase
nas condições que levam o Quicksort ao seu pior caso.

\section{Ordenação por inserção}
\label{sec:insertion}

O \textit{Insertion Sort}, ou ordenação por inserção, constrói o vetor ordenado
incrementalmente. A cada iteração $i$, o algoritmo assume que o prefixo
$v[0..i-1]$ já está ordenado e insere $v[i]$ na posição correta desse prefixo,
deslocando para a direita todos os elementos maiores que ele. É o método que,
segundo \citeonline{ziviani2018}, corresponde à forma como um jogador ordena as
cartas na mão.

Sua análise é direta. Seja $C(n)$ o número de comparações entre chaves. No
\textbf{melhor caso} --- vetor já ordenado --- cada elemento é comparado apenas com
seu predecessor imediato e nenhum deslocamento ocorre, resultando em

\begin{equation}
C_{\min}(n) = n - 1 = \Theta(n).
\end{equation}

No \textbf{pior caso} --- vetor em ordem inversa --- o elemento $v[i]$ precisa ser
comparado com todos os $i$ elementos anteriores, de modo que

\begin{equation}
C_{\max}(n) = \sum_{i=1}^{n-1} i = \frac{n(n-1)}{2} = \Theta(n^2).
\end{equation}

No \textbf{caso médio}, supondo permutações equiprováveis, cada elemento percorre
em média metade do prefixo, resultando em aproximadamente $n^2/4$ comparações, o
que também é $\Theta(n^2)$ \cite{cormen2022}.

Apesar da complexidade quadrática, o Insertion Sort possui três propriedades que o
tornam vantajoso para entradas pequenas, e que fundamentam sua utilização nas
versões híbridas deste trabalho:

\begin{enumerate}
  \item \textbf{Constantes multiplicativas muito baixas}: o laço interno consiste em
        uma comparação e uma atribuição, sem chamadas de função nem manipulação de
        pilha;
  \item \textbf{Adaptatividade}: em vetores quase ordenados o desempenho aproxima-se
        do linear, pois o número de deslocamentos é proporcional ao número de
        inversões da entrada;
  \item \textbf{Localidade de referência}: os acessos são sequenciais e
        contíguos, o que favorece a hierarquia de memória e a predição de desvios
        dos processadores modernos.
\end{enumerate}

\section{Quicksort}
\label{sec:quicksort}

O Quicksort \cite{hoare1961,hoare1962} é um algoritmo de divisão e conquista. Dado
um subvetor $v[esq..dir]$, escolhe-se um elemento como \textbf{pivô} e efetua-se a
\textbf{partição}: rearranja-se o subvetor de modo que todos os elementos menores ou
iguais ao pivô fiquem à sua esquerda e todos os maiores ou iguais, à sua direita.
Em seguida, o procedimento é aplicado recursivamente às duas partes. Diferentemente
do Mergesort, não há etapa de combinação: quando as chamadas recursivas retornam, o
vetor já está ordenado.

Este trabalho adota o esquema de partição de \citeonline{hoare1962} na forma
apresentada por \citeonline{ziviani2018}, em que dois índices caminham em sentidos
opostos até se cruzarem, trocando os elementos que se encontram no lado errado do
pivô. O pivô é tratado como \textbf{valor}, e não como posição; como esse valor é
sempre retirado do próprio subvetor, ele funciona como sentinela natural e garante a
terminação dos laços internos.

\subsection{Análise de complexidade}

O custo do Quicksort é determinado pelo \textbf{balanceamento das partições}. A
partição de um subvetor de tamanho $n$ custa $\Theta(n)$ comparações, e a recorrência
geral é

\begin{equation}
T(n) = T(k) + T(n - k - 1) + \Theta(n),
\end{equation}

onde $k$ é o tamanho da partição esquerda.

\paragraph{Melhor caso.} Quando o pivô é sempre a mediana exata, tem-se
$k \approx n/2$ e a recorrência torna-se $T(n) = 2T(n/2) + \Theta(n)$, cuja solução,
pelo Teorema Mestre, é

\begin{equation}
T_{\min}(n) = \Theta(n \log n).
\end{equation}

\paragraph{Caso médio.} Supondo que todas as permutações da entrada sejam
equiprováveis, o número esperado de comparações é

\begin{equation}
C_{\text{méd}}(n) \approx 2n \ln n \approx 1{,}386\, n \log_2 n,
\end{equation}

conforme demonstrado por \citeonline{hoare1962} e detalhado por
\citeonline{knuth1998}. O caso médio é, portanto, apenas cerca de 39\% pior que o
melhor caso --- resultado que explica a robustez prática do algoritmo.

\paragraph{Pior caso.} Quando o pivô é sistematicamente o menor ou o maior elemento
do subvetor, tem-se $k = 0$ a cada chamada, e a recorrência degenera para
$T(n) = T(n-1) + \Theta(n)$, cuja solução é

\begin{equation}
C_{\max}(n) = \sum_{i=1}^{n-1}(n - i) = \frac{n(n-1)}{2} \approx \frac{n^2}{2} = \Theta(n^2).
\label{eq:piorcaso}
\end{equation}

Além do custo quadrático em tempo, o pior caso produz uma cadeia de recursão de
profundidade $n$, elevando o consumo de memória de $O(\log n)$ para $O(n)$ e podendo
provocar estouro de pilha.

\subsection{Como o pior caso ocorre}
\label{sec:piorcaso}

O ponto central --- e frequentemente contraintuitivo --- é que o pior caso do
Quicksort \textbf{não é uma propriedade da entrada isoladamente, mas da interação
entre a entrada e a regra de escolha do pivô}. Não existe entrada intrinsecamente
``ruim''; existe entrada ruim \textit{para uma dada estratégia de pivô}.

Se a estratégia adotada é ``pivô = primeiro elemento'', um vetor já ordenado é o
pior caso: o menor elemento é escolhido como pivô a cada chamada, produzindo
partições de tamanhos $0$ e $n-1$. Esse é um cenário particularmente perigoso na
prática, pois vetores já ordenados ou quase ordenados são extremamente comuns em
sistemas reais --- exatamente a situação em que a implementação ingênua falha.

Se a estratégia é ``pivô = elemento central'', o vetor ordenado passa a ser o
\textbf{melhor} caso, pois o elemento central de um vetor ordenado é precisamente sua
mediana. Ainda assim, é possível construir entradas adversárias que forçam o
comportamento quadrático, como demonstrado por \citeonline{bentley1993}, que
apresentam um método sistemático para gerar entradas patológicas contra qualquer
regra determinística de escolha de pivô.

\section{Aprimoramentos clássicos}

\subsection{Corte para Insertion Sort}

A ideia de interromper a recursão em subvetores pequenos e finalizar a ordenação com
Insertion Sort foi proposta por \citeonline{singleton1969} e analisada em detalhe por
\citeonline{sedgewick1978}. A justificativa é que, embora o Insertion Sort seja
$\Theta(n^2)$, sua constante multiplicativa é muito menor que a do Quicksort; abaixo
de um limiar $M$, o termo quadrático com constante pequena é menor que o termo
$n\log n$ com constante grande.

\citeonline{sedgewick1978} sugere valores de $M$ entre 5 e 25, faixa reproduzida em
\citeonline{sedgewick2011}. Ressalte-se, entretanto, que tais valores foram obtidos
em hardware da década de 1970; um dos objetivos deste trabalho é justamente
determinar $M$ empiricamente na arquitetura atual.

\subsection{Mediana-de-três}

A escolha do pivô como mediana entre $v[esq]$, $v[meio]$ e $v[dir]$ também é atribuída
a \citeonline{singleton1969}. A técnica tem dois efeitos: melhora a estimativa da
mediana verdadeira do subvetor, reduzindo o desbalanceamento esperado das partições,
e elimina o comportamento quadrático em entradas ordenadas e inversamente ordenadas.

\citeonline{sedgewick1978} demonstra que a mediana-de-três reduz o número esperado de
comparações no caso médio de $2n\ln n$ para aproximadamente $\frac{12}{7} n \ln n$, uma
redução em torno de 14\%. É importante observar, contudo, que a técnica \textbf{não
elimina} o pior caso: ela apenas torna as entradas patológicas mais difíceis de
ocorrer por acaso. \citeonline{musser1997} demonstra esse ponto e propõe o
\textit{Introsort}, que monitora a profundidade da recursão e alterna para Heapsort
quando ela excede $2\log_2 n$, garantindo $O(n \log n)$ no pior caso. Essa é a
estratégia adotada pela função \texttt{std::sort} da biblioteca padrão do C++.

\section{Trabalhos recentes}

A pesquisa sobre Quicksort permanece ativa, com foco no aproveitamento das
características do hardware moderno.

\citeonline{wild2012} realizam a análise de caso médio do Quicksort de pivô duplo de
Yaroslavskiy, adotado no Java 7, mostrando que ele executa cerca de 5\% menos
comparações que a versão de pivô único --- resultado estendido por
\citeonline{kushagra2014} para múltiplos pivôs, que atribuem o ganho real à redução
do número de falhas de cache, e não à contagem de comparações.

\citeonline{edelkamp2019} propõem o \textit{BlockQuicksort}, que reorganiza a
partição em blocos a fim de eliminar desvios condicionais dependentes de dados,
obtendo ganhos de até 80\% em relação ao \texttt{std::sort} ao evitar erros de
predição de desvio. \citeonline{aumuller2021} simplificam a técnica usando o esquema
de partição de Lomuto. Na mesma linha, \citeonline{peters2021} apresenta o
\textit{pattern-defeating quicksort} (pdqsort), que combina detecção de padrões na
entrada com partição sem desvios, hoje adotado pela biblioteca padrão da linguagem
Rust. \citeonline{axtmann2022} estendem essas ideias para o contexto paralelo de
memória compartilhada.

O ponto comum a esses trabalhos é relevante para a análise deste trabalho: em
hardware moderno, \textbf{o número de comparações deixou de ser um bom preditor
isolado do tempo de execução}. Efeitos de cache e de predição de desvios respondem
por parcela significativa do custo real, o que será observado experimentalmente no
Capítulo \ref{cap:resultados}.
